% Chapter 1 -- Using Initech 123 (annex Sec 6, Reference Chapter 1: the
% screen, the indicators, the keys, the menus, entering data, ranges,
% formulas, recalculation, files).
% Sources (Law 1): os/i123/app/face_a.c (the screen layout: control panel
% lines 1-3, the column and row borders, the status line with CALC / CIRC;
% MODE_TEXT; key handling in key() and command_key(); menu_select; the
% prompts and notices; commit(): a value that does not parse drops into
% EDIT; request_quit: the exit guard), os/i123/core/fmt.c (label prefixes,
% i123_parse_value: the trailing %, the E exponent), os/i123/core/formula.c
% and spec/i123/wk1_formula.h (operators and their precedence),
% os/i123/core/recalc.c (natural / columnwise / rowwise order, automatic
% and manual recalculation, iteration, circular references),
% os/i123/core/range.c (range text: A1..B2, A1.B2, names).  The key table
% registers every key the probe finds acting (harness/pubs/i123_face.c)
% and, DRAFT, the Release 2.2 keys this version does not act on.
% GATES: test-i123-wk1-roundtrip test-i123-formula test-i123-recalc-order test-i123-range-menu test-i123-quit test-flair-i123 test-pubs-i123 test-pubs-i123-examples test-pubs-screens
\Gates{test-i123-wk1-roundtrip, test-i123-formula, test-i123-recalc-order,
  test-i123-range-menu, test-i123-quit, test-flair-i123, test-pubs-i123,
  test-pubs-i123-examples, test-pubs-screens}
\chapter{Using Initech 123}

This chapter describes the parts of the Initech 123 screen, the keys you
use, the way you choose commands from the menus, and the kinds of
information you can enter in a worksheet.  Read it before you use the rest
of this book: the command chapters assume that you know the terms it
introduces.

\section{Starting and Ending Initech 123}
\ix{starting Initech 123}
Initech 123 is a FLAIR application.  Its program file, \file{123.EXE}, is
in the folder \file{APPS} of the InitechOS data disk, together with
\file{123.MNU}, the file that holds its menus.  To start Initech 123, open
the \file{APPS} folder and double-click the \file{123.EXE} icon.  Initech
123 opens its window, titled \texttt{Initech 123}, with a new, empty
worksheet.

\ix{ending Initech 123}
To end Initech 123, select \slashcmd{/Quit Yes}, or choose \textbf{Quit}
from the \textbf{File} menu of the menu bar (keyboard equivalent
\key{Ctrl-Q}), or click the close box of the window.  Each of these first
checks the worksheet: if you have made a change since you last saved or
retrieved it, or if you are in the middle of typing an entry, Initech 123
asks whether to save it before it ends (see Chapter~9, ``The Quit
Command'').

\begin{Note}
Initech 123 needs \file{123.MNU} in the \file{APPS} folder.  If the file is
missing or damaged, Initech 123 does not start.
\end{Note}

\section{The Initech 123 Screen}
\ix{screen}
Figure~\ref{fig:screen} shows the Initech 123 window with a small
worksheet in it.

\FlairShot[fig:screen]{screen}{The Initech 123 window}

The window has four areas.
\begin{itemize}
\item The \term{control panel}, the top three lines of the window.  The
  first line shows the address of the cell the cell pointer is in, a colon
  and the cell's contents, and, at the right, the mode indicator.  The
  second line shows the entry you are typing, the menu you are using, or a
  prompt.  The third line describes the highlighted menu item.
\item The \term{worksheet frame}: the column letters across the top of the
  worksheet area and the row numbers down its left side, both in reverse
  video.
\item The \term{worksheet area}, which shows as many whole columns as fit
  in the window (eight columns of the standard width) and nineteen rows.
\item The \term{status line}, the bottom line of the window.  It shows
  messages, such as the name of the file you last saved, and, at the right,
  the status indicators.
\end{itemize}

\subsection{The Worksheet}
\ix{worksheet!size}
A \term{worksheet} has 256 columns, lettered A through Z, AA through AZ, and
so on to IV, and 8,192 rows, numbered 1 through 8,192.  The intersection of a
column and a row is a \term{cell}.  A cell is named by its column letter
and row number, its \term{address}: \texttt{A1} is the upper left cell of
the worksheet, \texttt{IV8192} the lower right.

The \term{cell pointer} is the highlighted bar that marks the
\term{current cell}, the cell that receives what you type.  It is as wide
as the column it is in.

\subsection{Mode Indicators}
\ix{mode indicators}
The \term{mode indicator} at the right of the first line of the control
panel tells you what Initech 123 is doing and what it expects from you.

\begin{LotTable}{p{24mm}p{104mm}}
\LotHead{\textbf{Indicator} & \textbf{Meaning}}
\ind{READY} & Initech 123 is waiting for you to make an entry or to select a
  command.\\
\ind{LABEL} & You are entering a label.\\
\ind{VALUE} & You are entering a value: a number or a formula.\\
\ind{EDIT} & Initech 123 found an error in a value you entered and is
  waiting for you to correct it.\\
\ind{MENU} & A menu is displayed in the control panel.\\
\ind{FILES} & Initech 123 is asking for the name of a file.\\
\ind{RANGE} & Initech 123 is asking for a range, an address, a range name
  or a number.\\
\ind{ERROR} & Initech 123 is displaying a message on the second line of
  the control panel.  Press \kname{ESC} to return to the
  worksheet.\LotEnd
\end{LotTable}

\subsection{Status Indicators}
\ix{status indicators}
The status indicators appear in reverse video at the right end of the
status line.

\begin{LotTable}{p{24mm}p{104mm}}
\LotHead{\textbf{Indicator} & \textbf{Meaning}}
\ind{CALC} & Recalculation is manual, and an entry has changed since the
  worksheet was last recalculated.  Press \fkey{Calc}{F9} to recalculate.\\
\ind{CIRC} & The worksheet contains a circular reference: a formula that
  refers to itself, directly or through other formulas.\LotEnd
\end{LotTable}

\section{The Keyboard}
\ix{keys}
This section lists the keys Initech 123 uses.  A key marked \DraftSample\
has a function in the Initech 123 design that this version does not yet
carry out; pressing it has no effect.

\subsection{Pointer-Movement Keys}
\ix{pointer-movement keys}
In \ind{READY} mode the pointer-movement keys move the cell pointer.  When
the cell pointer reaches the edge of the window, the worksheet scrolls.
While you are typing a label or a value, a pointer-movement key first
completes the entry, as \kname{ENTER} does, and then moves the cell
pointer.  At a range prompt, the arrow keys move the cell pointer to the
cell you want to name.

\begin{KeyTable}{Key}{Name}
\KeyRow{UP}{Up arrow}{Moves the cell pointer up one row.}
\KeyRow{DOWN}{Down arrow}{Moves the cell pointer down one row.}
\KeyRow{LEFT}{Left arrow}{Moves the cell pointer left one column.}
\KeyRow{RIGHT}{Right arrow}{Moves the cell pointer right one column.}
\KeyRow{HOME}{Home}{Moves the cell pointer to cell A1.}
\KeyRow[not available in this version]{PGUP}{PgUp}{Moves the cell pointer
  up one window.}
\KeyRow[not available in this version]{PGDN}{PgDn}{Moves the cell pointer
  down one window.}
\KeyRow[not available in this version]{TAB}{Tab}{Moves the cell pointer
  right one window.}
\KeyRow[not available in this version]{END}{End}{Used with an arrow key,
  moves the cell pointer to the end of a block of entries.}
\end{KeyTable}

\subsection{Function Keys}
\ix{function keys}
Each function key that Initech 123 uses has a name, printed on the
keyboard template.  This book refers to a function key by its name with
the key in parentheses: \fkey{Goto}{F5}.

\begin{KeyTable}{Key}{Name}
\KeyRow[not available in this version]{F1}{HELP}{Displays a help screen.}
\KeyRow{F2}{EDIT}{In \ind{READY} mode, places the contents of the current
  cell on the second line of the control panel so that you can change
  them.}
\KeyRow[not available in this version]{F3}{NAME}{Lists the range names.}
\KeyRow[not available in this version]{F4}{ABS}{Makes a cell reference
  absolute.}
\KeyRow{F5}{GOTO}{In \ind{READY} mode, asks for a cell address or a range
  name and moves the cell pointer there.}
\KeyRow[not available in this version]{F6}{WINDOW}{Moves the cell pointer
  to the other window.}
\KeyRow[not available in this version]{F7}{QUERY}{Repeats the last
  query.}
\KeyRow[not available in this version]{F8}{TABLE}{Repeats the last data
  table.}
\KeyRow{F9}{CALC}{In \ind{READY} mode, recalculates every formula in the
  worksheet.}
\KeyRow[not available in this version]{F10}{GRAPH}{Displays the current
  graph.}
\end{KeyTable}

\subsection{Special Keys}
\ix{special keys}
\begin{KeyTable}{Key}{Name}
\KeyRow{SLASH}{/}{In \ind{READY} mode, displays the main menu.}
\KeyRow{LESS-THAN}{<}{The same as \kname{SLASH}.}
\KeyRow{ENTER}{Enter}{Completes an entry, a menu choice or a prompt.}
\KeyRow{ESC}{Esc}{Cancels the entry you are typing; in a menu, returns to
  the previous menu; at a prompt, returns to the menu; after a message,
  returns to the worksheet.}
\KeyRow{BACKSPACE}{Backspace}{Erases the character to the left of the
  cursor in an entry or a prompt.}
\KeyRow{CTRL-BREAK}{Ctrl-Break}{Cancels a menu or a prompt and returns to
  \ind{READY} mode.}
\KeyRow[not available in this version]{DEL}{Del}{In \ind{EDIT} mode,
  erases the character at the cursor.}
\KeyRow[not available in this version]{INS}{Ins}{In \ind{EDIT} mode,
  switches between inserting and overtyping.}
\end{KeyTable}

\begin{Note}
On a keyboard without a \key{Break} key, \key{Ctrl-C} does the same as
\kname{CTRL-BREAK}.
\end{Note}

\section{Using the Menus}
\ix{menus}
You give Initech 123 commands by selecting them from menus in the control
panel.  Press \kname{SLASH} in \ind{READY} mode to display the
\term{main menu}.  The mode indicator changes to \ind{MENU}, the second line
of the control panel lists the commands of the menu with the first one
highlighted by the \term{menu pointer}, and the third line describes the
highlighted command (Figure~\ref{fig:menumain}).  If the highlighted
command opens another menu, the third line lists the commands of that
menu.

\FlairShot[fig:menumain]{menumain}{The main menu}

To select a command, do one of the following:
\begin{itemize}
\item Move the menu pointer to the command with \kname{RIGHT} or
  \kname{LEFT} (\kname{DOWN} and \kname{UP} work the same way; \kname{HOME}
  goes to the first command) and press \kname{ENTER}.
\item Type the first character of the command, in capital or small
  letters.  If two commands of a menu begin with the same character, typing
  it selects the first of them.
\end{itemize}
Most commands open a further menu; you continue selecting until you reach
the command you want.  Figure~\ref{fig:menurf} shows the menu that appears
when you select Range and then Format.  This book names a command by the
words you select, after a slash: \slashcmd{/Range Format Currency} means:
press \kname{SLASH}, select Range, then Format, then Currency.

\FlairShot[fig:menurf]{menurf}{The /Range Format menu}

To back out of a menu, press \kname{ESC}: Initech 123 returns to the
previous menu, with the command you selected from it still highlighted.
\kname{ESC} at the main menu returns to \ind{READY} mode.  To leave the
menus at once from any level, press \kname{CTRL-BREAK}.

\subsection{Commands Not Available in This Version}
\ix{Command not implemented}
The menus of Initech 123 are complete: every command of the Initech 123
design appears in them, in its place.  Some of those commands are not yet
available in this version.  When you select one, the mode indicator changes
to \ind{ERROR}, the control panel displays \texttt{Command not implemented},
and nothing in the worksheet changes (Figure~\ref{fig:refused}).  Press
\kname{ESC} to return to the worksheet.

\FlairShot[fig:refused]{refused}{A command that is not available in this
version}

In this book such commands are described within passages marked \DraftSample.

\subsection{Prompts}
\ix{prompts}
Many commands ask you for information: a range, a file name, a number.  The
question, the \term{prompt}, appears on the second line of the control
panel.  Type the answer and press \kname{ENTER}.  While you type,
\kname{BACKSPACE} erases the last character.  \kname{ESC} returns you to
the menu you came from; \kname{CTRL-BREAK} returns to \ind{READY} mode.

When a prompt asks for a range, the control panel shows the address of the
current cell after the prompt.  Press \kname{ENTER} at once to accept it,
type a range, or move the cell pointer with the arrow keys to another cell
and press \kname{ENTER}.

If Initech 123 cannot accept the answer, it displays a message on the
second line of the control panel, such as \texttt{Invalid range} or
\texttt{Illegal number}, and the mode indicator changes to \ind{ERROR}.
Press \kname{ESC} to return to the worksheet, and select the command again.

\section{Entering Data}
\ix{entries}
To make an entry, move the cell pointer to the cell, type the entry and
press \kname{ENTER}, or press one of the pointer-movement keys to complete
the entry and move the cell pointer in one step.  While you type, the entry
appears on the second line of the control panel, followed by the cursor.
Press \kname{BACKSPACE} to erase the last character, or \kname{ESC} to
abandon the entry.

An entry you type can be up to 74 characters long; further characters are
not accepted.  Initech 123 decides what kind of entry you are making from
its first character:
\begin{itemize}
\item an entry that begins with a digit, or with one of the characters
  \texttt{+ - . ( @ \# \$}, is a \term{value}; the mode indicator shows
  \ind{VALUE};
\item any other entry is a \term{label}; the mode indicator shows
  \ind{LABEL} (Figure~\ref{fig:labelmode}).
\end{itemize}

\FlairShot[fig:labelmode]{labelmode}{Typing a label}

\subsection{Labels}
\ix{labels}
A \term{label} is text.  Initech 123 stores a label with a
\term{label prefix}, a character that says how to align the label in its
cell.  The prefix is not displayed in the cell; it appears before the label
on the first line of the control panel.

\begin{LotTable}{p{24mm}p{104mm}}
\LotHead{\textbf{Prefix} & \textbf{Alignment}}
\texttt{'} (apostrophe) & Left-aligned\\
\texttt{"} (quotation mark) & Right-aligned\\
\texttt{\textasciicircum} (caret) & Centered\\
\texttt{\textbackslash} (backslash) & Repeated to fill the cell\LotEnd
\end{LotTable}

If you do not type a prefix, Initech 123 adds the default prefix, the
apostrophe unless you have changed it with
\slashcmd{/Worksheet Global Label-Prefix}.  To enter a label that begins
with a digit or with one of the value characters, such as a year or a part
number, type a prefix first.  These entries, and the cells as they then
appear:

\LotExample{examples/labels}

A left-aligned label longer than its column continues into the cells to its
right as long as they are empty.

\subsection{Values}
\ix{values}
A \term{value} is a number or a formula.  Type a number with or without a
sign and a decimal point.  You can type a number in scientific notation,
with \texttt{E} and a power of ten (\type{1.5E3} is 1500), and you can end a
number with a percent sign to divide it by 100 (\type{15.5\%} is 0.155).
Do not type commas, spaces or dollar signs in a number: the cell's format
(Chapter~10, \slashcmd{/Range Format}) decides how the number is
displayed.

\LotExample{examples/values}

If Initech 123 cannot accept a value, because it contains a character a
number or formula cannot contain, it switches to \ind{EDIT} mode and leaves
the entry on the second line of the control panel so that you can correct
it.  Press \kname{ENTER} again when it is correct, or \kname{ESC} to
abandon it.

\subsection{Editing an Entry}
\ix{editing}
To change the entry in a cell, move the cell pointer to it and press
\fkey{Edit}{F2}.  Initech 123 places the contents of the cell, with its
label prefix if it is a label, on the second line of the control panel,
and the mode indicator shows \ind{LABEL} or \ind{VALUE}
(Figure~\ref{fig:editmode}).  Erase from the end with \kname{BACKSPACE},
type the new characters, and press \kname{ENTER}.  An arrow key completes
the entry as \kname{ENTER} does and then moves the cell pointer;
\kname{ESC} abandons the change and leaves the cell as it was.

\FlairShot[fig:editmode]{editmode}{A formula placed on the control panel by
EDIT (F2)}

\subsection{Erasing an Entry}
To erase the contents of cells, use \slashcmd{/Range Erase}
(Chapter~10).  To erase the whole worksheet, use
\slashcmd{/Worksheet Erase Yes} (Chapter~12).

\section{Ranges}
\ix{ranges}
A \term{range} is a rectangle of cells: a single cell, part of a row or a
column, or a block of rows and columns.  You name a range by the addresses
of two opposite corners separated by two periods, \type{A1..C5}; one period
is also accepted, \type{A1.C5}.  The corners can be given in any order.  A
single cell is named by its address alone.

A \term{range name} is a name of up to 15 characters that you give a range
with \slashcmd{/Range Name Create}.  You can then type the name wherever
Initech 123 asks for a range, and in formulas.

\section{Formulas}
\ix{formulas}
A \term{formula} calculates a value from numbers, cell references, range
references and @functions.  A formula must begin with one of the value
characters: to begin a formula with a cell address, type a plus sign
first, as in \type{+B2*B3}.  When you complete the entry, the cell shows
the result, and the first line of the control panel shows the formula.

\subsection{Operators}
\ix{operators}
The operators combine the parts of a formula.  Initech 123 performs the
operations in the order of the table below, the highest precedence first;
operations of the same precedence are performed from left to right.
Parentheses change the order: what is inside them is calculated first.

\begin{LotTable}{p{24mm}p{74mm}p{26mm}}
\LotHead{\textbf{Operator} & \textbf{Operation} & \textbf{Precedence}}
\texttt{\textasciicircum} & Exponentiation & 7 (highest)\\
\texttt{-} \texttt{+} & Negative, positive & 6\\
\texttt{*} \texttt{/} & Multiplication, division & 5\\
\texttt{+} \texttt{-} & Addition, subtraction & 4\\
\texttt{= <> < > <= >=} & Comparison (true is 1, false is 0) & 3\\
\texttt{\#NOT\#} & Logical NOT & 2\\
\texttt{\#AND\# \#OR\#} & Logical AND, logical OR & 1\\
\texttt{\&} & String concatenation & 1 (lowest)\LotEnd
\end{LotTable}

These formulas and their results show the order of operations:

\LotExample{examples/operators}

\subsection{Cell and Range References}
\ix{references}
A formula can refer to a cell by its address and to a range by its corners
or its name.  A \term{relative reference}, such as \texttt{A1}, is stored as
a distance from the formula's own cell; an \term{absolute reference}, such
as \texttt{\$A\$1}, is stored as a fixed address.  A \term{mixed
reference}, such as \texttt{\$A1} or \texttt{A\$1}, fixes only the column or
only the row.  The difference matters when you copy a formula: the copy of
a relative reference refers to the cell at the same distance from the copy,
the copy of an absolute reference to the same cell.  Here B1 and C1 are
copied to B2 and C2:

\LotExample{examples/refs}

\subsection{@Functions}
An \term{@function} is a built-in formula, such as \fn{@SUM} or
\fn{@ROUND}, that you use as part of a formula.  Chapter~13 describes every
@function.

\subsection{Formulas with Errors}
If a formula cannot be calculated, for instance because it divides by zero
or refers to a label where it needs a number, the cell shows \texttt{ERR}.
Every formula that depends on such a cell also shows \texttt{ERR}, until the
cause is corrected.

\section{Recalculation}
\ix{recalculation}
Initech 123 recalculates the formulas of the worksheet whenever you change
an entry, in \term{natural order}: before it calculates a formula, it
calculates every formula the formula depends on.  You can change both when
and in what order it recalculates, with
\slashcmd{/Worksheet Global Recalculation} (Chapter~12).

When recalculation is manual, Initech 123 does not recalculate after an
entry; the \ind{CALC} indicator appears on the status line instead
(Figure~\ref{fig:calcind}).  Press \fkey{Calc}{F9} to recalculate.

\FlairShot[fig:calcind]{calcind}{Manual recalculation: B1 has changed, and
the total in B3 waits for CALC (F9)}

\subsection{Circular References}
\ix{circular reference}
A \term{circular reference} is a formula that depends on its own value,
directly or through other formulas.  When the worksheet contains one, the
\ind{CIRC} indicator appears on the status line (Figure~\ref{fig:circind}).
Each recalculation then calculates the formulas once (or as many times as
\slashcmd{/Worksheet Global Recalculation Iteration} specifies), so the
values of the formulas in the circle change each time you recalculate.

\FlairShot[fig:circind]{circind}{A circular reference: B2 adds itself and
B3, and B3 adds B2}

\section{Worksheet Files}
\ix{files}
A worksheet exists in memory until you save it in a \term{worksheet file}
with \slashcmd{/File Save}.  Initech 123 saves and retrieves worksheet
files in the \file{.WK1} format.  A file name has up to eight characters,
optionally followed by a period and an extension of up to three characters;
when you do not type an extension, Initech 123 adds \file{.WK1}.  The data
disk is drive \file{A:}.  Chapter~5 describes the File commands.
